\scanpage{39}
\mspage{\handwritten{7}}
% archive sheet number: 38
$G_{\handwritten{\eta}}$ est séparé sur $\handwritten{\kappa}(\handwritten{\eta})$ (\deleted{Exp.} $\mathrm{VI}_A$ 0.2), d'après (10.2,10.10 a),EGA IV 8.8.2 et 8.10.5), il existe un ouvert\inserted{affine} $S'$ de $S$ tel que $G'$ soit un $S'$-groupe de \deleted{type}\inserted{présentation} fini, plat et séparé sur $S'$, et que $X'$, $H'$ et $K'$ soient plats et de présentation finie sur $S'$. Si on suppose $H_{\handwritten{\eta}}$ (resp\unsure{;|.} $K_{\handwritten{\eta}}$) fermé dans $G_{\handwritten{\eta}}$, on peut choisir $S'$ de sorte que $H'$ (resp. $K'$) soit fermé dans $G'$. D'après (EGA IV 11.3.15), quitte à remplacer $S'$ par un ouvert affine de $S'$, on peut supposer que $G'$, $H'$, $K'$ et $X'$ sont essentiellement libres sur $S'$ (au sens de \deleted{Exp.} VIII (6.1)). Il résulte alors de (\deleted{Exp.} VIII \handwritten{\unsure{5|6}}.5 b) et e)\textsuperscript{(*)}) que, sous les hypothèses de l'énoncé, les foncteurs considérés sont représentables par des sous-préschémas fermés (donc de présentation finie sur $S'$) de $G'$.

\uline{Proposition 10.12.- Soient $S$ un préschéma intègre, $G$ un \deleted{$S$-préschéma} $S$-groupe de présentation finie, $A$ et $B$ deux sous-préschémas en groupes de présentation finie de $G$, à fibre générique lisse, tels que \deleted{ou bien $A$ soit à fibres connexes}\inserted{$A$ soit à fibre générique connexe}, \deleted{ou bien bien}\inserted{(resp.} $A$ et $B$ soient invariants dans $G$). Alors il existe un ouvert non vide $S'$ de $S$ et un sous-préschéma en groupes\inserted{\handwritten{fermé, $D'$}} de présentation finie, à fibres lisses \deleted{et fermé $D'$} de $G'=G\handwritten{|}S'$\handwritten{,} tel que $D'$ soit \deleted{\unsure{ce|co}} à fibres connexes (resp. soit invariant dans $G'$), et que $D'$ représente le faisceau associé, pour la topologie fppf ou fpqc, au préfaisceau en groupes des commutateurs de $A$ et $B$ dans $G$. En particulier, pour tout $s\handwritten{\in}S'$, on a $D'_s=(A_s,B_s)$ avec les notations de (7.2 (vii)).}

Soit $\handwritten{\eta}$ le point générique de $S$; posons $D'_{\handwritten{\eta}}=(A_{\handwritten{\eta}},B_{\handwritten{\eta}})$. D'après (7.8) (resp. 7.3 v)), $D'_{\handwritten{\eta}}$ est connexe (resp. invariant dans $G_{\handwritten{\eta}}$). D'après (10.1 et 10.10 a)), il existe un ouvert non vide $S'$ de $S$ et un sous-$S'$-préschéma en groupes $D'$ de présentation finie et fermé dans $G'$, à fibres séparables (EGA IV 9.7.7 et 9.3.3), ayant $D'_{\handwritten{\eta}}$ pour fibre au point $\handwritten{\eta}$, et tel que $D'$ soit à fibres connexes (i.e. géométriquement connexes (\deleted{Exp.} $\mathrm{VI}_A$ 2.1.1))(resp. soit invariant dans $G'$)(EGA IV 9.7.7 et 9.3.3\inserted{\handwritten{,}}(resp. (10.4)). De plus, nous avons vu, au cours de la démonstration de (7.8), qu'il existe un entier $n$ tel que $\handwritten{\gamma}_{\handwritten{\eta}}^n((A_{\handwritten{\eta}}\times_{\handwritten{\kappa}(\handwritten{\eta})}B_{\handwritten{\eta}})^n)=D'$, où $\handwritten{\gamma}_{\handwritten{\eta}}:A_{\handwritten{\eta}}\times_{\handwritten{\kappa}(\handwritten{\eta})}B_{\handwritten{\eta}}\handwritten{\longrightarrow}G_{\handwritten{\eta}}$

\par\noindent\handwritten{\rule{\linewidth}{0.3pt}}\par
\handwritten{(*) Voir note page 51 \textcircled{40}}\par
